\chapter{Առաջարկվող MVP (նվազագույն կենսունակ տարբերակ)}

Բակալավրիատի ծավալը իրատեսական պահելու համար առաջին տարբերակը պետք է պարունակի.

\textbf{Ապարատային մաս.}
\begin{itemize}[leftmargin=1.4em]
  \item մեկ AUP ZU3 հարթակ,
  \item 3 էմուլացված սենսորի տիպ՝ Temperature, Pressure, Vibration,
  \item վերաօգտագործելի գեներատորներ՝ Ramp, Periodic, LFSR/Noise, Anomaly injection,
  \item FIFO/BRAM բուֆերացում,
  \item մոտ 3 processing IP՝ Moving Average, Min/Max, Threshold/Event Detection,
  \item AXI4-Lite կոնֆիգուրացիա,
  \item PL$\rightarrow$PS տվյալների փոխանցում,
  \item Ethernet հաղորդակցություն։
\end{itemize}

\textbf{Ծրագրային մաս.}
\begin{itemize}[leftmargin=1.4em]
  \item մեկ սերվերային համակարգիչ (Backend/API, տվյալների բազա),
  \item Web dashboard՝ փորձի կոնֆիգուրացիայով,
  \item տվյալների վիզուալիզացիա (live և պատմական),
  \item փորձերի համեմատում,
  \item հիմնական վիճակագրական վերլուծություն,
  \item մեկ մոդելավորման / անոմալիաների հայտնաբերման կիրառություն։
\end{itemize}
